\documentclass[12pt,oneside]{book}

\usepackage[left=2cm,right=2cm,top=2.5cm,bottom=2.5cm]{geometry}
\usepackage[unicode=true,colorlinks=true,linkcolor=blue]{hyperref}

\usepackage{amsmath}
\usepackage{amsfonts}
\usepackage{amssymb}
\usepackage{amsthm}
\usepackage{mathtools}
\usepackage{mathrsfs}
\usepackage{cases}
\usepackage{pgf,tikz}
\usepackage{pgfplots}
\usetikzlibrary{shapes}
\usetikzlibrary{shapes.arrows}
\usetikzlibrary{arrows.meta}
\usetikzlibrary{calc}
\usetikzlibrary{math}
\usetikzlibrary{decorations,
    decorations.footprints,
    decorations.fractals,
    decorations.markings,
    decorations.pathmorphing,
    decorations.pathreplacing,
    decorations.shapes,
    decorations.shapes,
    decorations.text}
\usetikzlibrary{positioning}
\usetikzlibrary{angles}
\usetikzlibrary{matrix}
\usepackage{tikz-cd}

\usepackage{soul}
\usepackage{fancyhdr}
\usepackage{xcolor}
\usepackage{titlesec}
\usepackage{indentfirst}
\usepackage{chngcntr}
\usepackage{caption}
\usepackage{subcaption}
\usepackage{booktabs}
\usepackage[inline]{enumitem}
\usepackage{setspace}

\setstretch{1.4142}
\setcounter{chapter}{0}
\setcounter{section}{0}

\captionsetup{labelfont={bf},labelsep=period}
\counterwithin{figure}{chapter}
\counterwithin{table}{chapter}
\counterwithout{section}{chapter}
\renewcommand{\thesubsection}{\thesection\Alph{subsection}}

\theoremstyle{definition}

\newtheorem{exercise}{Exercise}
\counterwithin{exercise}{section}
\newenvironment{sqcases}{%
    \matrix@check\sqcases\env@sqcases
}{%
    \endarray\right.%
}
\def\env@sqcases{%
\let\@ifnextchar\new@ifnextchar
\left\lbrack{}
\def\arraystretch{1.2}%
\array{@{}l@{\quad}l@{}}%
}
\newtheorem{innercustomgeneric}{\customgenericname}
\providecommand{\customgenericname}{}
\newcommand{\newcustomtheorem}[2]{%
  \newenvironment{#1}[1]
  {%
   \renewcommand\customgenericname{#2}%
   \renewcommand\theinnercustomgeneric{##1}%
   \innercustomgeneric%
  }
  {\endinnercustomgeneric}
}
\newcustomtheorem{problem}{Problem}
\newcustomtheorem{proposition}{Proposition}

\newcommand{\tr}[1]{\left({#1}\right)}
\newcommand{\card}[1]{\left\vert{#1}\right\vert}
\newcommand{\rank}{\text{rank}}
\newcommand{\abs}[1]{\left\vert{#1}\right\vert}
\newcommand{\norm}[1]{\left\Vert{#1}\right\Vert}
\newcommand{\anglebracket}[1]{\left\langle{#1}\right\rangle}
\newcommand{\floor}[1]{\left\lfloor{#1}\right\rfloor}
\newcommand{\ceil}[1]{\left\lceil{#1}\right\rceil}
\newcommand{\openinterval}[1]{\left\rbrack{#1}\right\lbrack}
\newcommand{\closedinterval}[1]{\left\lbrack{#1}\right\rbrack}
\newcommand{\halfopenleft}[1]{\left\rbrack{#1}\right\rbrack}
\newcommand{\halfopenright}[1]{\left\lbrack{#1}\right\lbrack}
\newcommand{\set}[1]{\left\{{#1}\right\}}
\newcommand{\tuple}[1]{\left({#1}\right)}
\newenvironment{section*}[1]{% \begin{section*}{section title}....\end{section*}
  \section*{#1}
  \renewcommand\thesection{\thechapter.S}
  \setcounter{exercise}{0}}{}

\title{Stephen Willard's General Topology: Notes and Problems}
\author{Ngo Quang Duong}
\date{\today}

\begin{document}

\maketitle

\tableofcontents

\chapter{Set Theory and Metric Spaces}

\section{Set theory}

\subsection{Russell's Paradox}

The phenomenon to be presented here was first exhibited by Russell in 1901, and consequently is known as Russell's Paradox.

Suppose we allow as sets things \(A\) for which \( A \in A \). Let \( P \) be the set of all sets. Then \( P \) can be divided into two nonempty subsets, \( P_{1} = \left\{ A \in P \mid A \notin A \right\} \) and \( P_{2} = \left\{ A \in P \mid A \in A \right\} \). Show that this results in the contradiction: \( P_{1} \in P_{1} \iff P_{1} \notin P_{1} \). Does our (na{\i}ve) restriction on sets given in 1.1 eliminate the contradiction?

\hrulefill%

If \( P_{1} \in P_{1} \) then \( P_{1} \notin P_{1} \), according to the definition of \( P_{1} \). Conversely, if \( P_{1} \notin P_{1} \) then \( P_{1} \in P_{1} \).

Our na{\i}ve restriction on sets (no aggregation shall be a set which would be an element of itself) does not eliminate the contradiction.

\subsection*{1B. De Morgan's laws and the distributive laws}

\begin{enumerate}[itemsep=0pt]
    \item \( A - \left( \bigcap_{\lambda \in \Lambda} B_{\lambda} = \bigcup_{\lambda \in \Lambda} (A - B_{\lambda}) \right) \)
    \item \( A - \left( \bigcap_{\lambda \in \Lambda} B_{\lambda} = \bigcup_{\lambda \in \Lambda} (A - B_{\lambda}) \right) \)
    \item If \( A_{nm} \) is a subset of \( A \) for \( n = 1, 2, \ldots \) and \( m = 1, 2, \ldots \) is it necessarily true that
    \[
        \bigcup_{n=1}^{\infty} \left\lbrack \bigcap_{m=1}^{\infty} A_{nm} \right\rbrack = \bigcap_{n=1}^{\infty} \left\lbrack \bigcup_{m=1}^{\infty} A_{nm} \right\rbrack?
    \]
\end{enumerate}

\subsection*{1C. Ordered pairs}

Show that, if \( (x_{1}, x_{2}) \) is defined to be \( \left\{ \left\{ x_{1} \right\}, \left\{ x_{1}, x_{2} \right\} \right\} \), then \( (x_{1}, x_{2}) = (y_{1}, y_{2}) \) iff \( x_{1} = y_{1} \) and \( x_{2} = y_{2} \).

\subsection*{1D. Cartesian products}

\section{Metric spaces}

\chapter{Topological Spaces}

\section{Fundamental concepts}

\section{Neighborhoods}

\section{Bases and subbases}

\chapter{New Spaces from Old}

\section{Subspaces}

\section{Continuous functions}

\section{Product spaces; weak topologies}

\section{Quotient spaces}

\chapter{Convergence}

\section{Inadequacy of sequences}

\section{Nets}

\section{Filters}

\chapter{Separation and Countability}

\chapter{Compactness}

\chapter{Metrizable Spaces}

\chapter{Connectedness}

\chapter{Uniform Spaces}

\chapter{Function Spaces}


\end{document}
